\chapter{Пересмотр приоритета: от одиночной нити к многопроходному узлу ткани}
\label{ch:v6-single-thread-fabric-reprioritization-gate}

\section{Исправленная исходная постановка}

Архивная формулировка в тексте \texttt{print2.md} не вводила независимое
множество фундаментальных струн. Она предполагала
одну непрерывную замкнутую нить конечного сечения, которая многократно
проходит через локальные области и после крупномасштабного усреднения
выглядит как ткань.

Нулевая толщина в этой постановке ведёт к бесконечной длине, необходимой
для заполнения положительного объёма. При ненулевом натяжении это означает
бесконечную энергию. Конечное сечение превращает локальные проходы одной
нити в задачу упаковки и контактной совместимости.

Строгая ветвь проекта не вывела микроскопическую нить и её планковское
сечение. Она, однако, независимо получила поле одноосного порядка $Q$,
тетраэдрический порядок $T$, устойчивую вихревую нить конечной ширины,
топологический дефицит замыкания и точные переносные нулевые моды.

Континуальные теории длинных связанных цепей подтверждают, что локальные
ориентационные моменты одной цепи подчинены дополнительным законам
сохранения, отсутствующим у набора независимых частиц
\cite{single-thread-moment-svensek2013,single-thread-moment-svensek2016}.
Тензор третьего ранга действительно способен хранить тетраэдрический
порядок при нулевых низших моментах
\cite{single-thread-moment-radzihovsky2000,single-thread-moment-lubensky2002}.

\section{Что именно меняет приоритет}

No-go спокойного кольца сохраняется: основное состояние одной эффективной
вихревой трубки сжимается под действием натяжения. Но этот результат не
проверяет частицу как локальный узел нескольких соседних проходов одной
глобальной нити.

Поэтому расчёт кривизного ответа одиночной трубки понижается до вторичного
вопроса об упругости ребра ткани. Первичным становится мост
\begin{equation}
 \text{одна глобальная нить}
 \longrightarrow
 \text{локальные проходы}
 \longrightarrow
 (Q,T,B)
 \longrightarrow
 \text{переносимый узел ткани}.
\end{equation}

Это не ретроспективное объявление ранней метафоры доказанной. Новая ветвь
обязана показать, что уже выведенные поля являются допустимыми моментами
одной связной линии, что локальные сшивки образуют один глобальный цикл и
что дефект сшивки имеет конечную энергию и сохраняемый класс.

\section{Предрегистрированная последовательность}

Проверки выполняются в следующем порядке:
\begin{enumerate}
 \item совместная моментная реализация $Q$ и $T$ локальными проходами;
 \item глобальная однокомпонентная сшивка локальных проходов;
 \item конечное сечение, упаковочная энергия и абсолютный масштаб;
 \item топологический класс многопроходного узла;
 \item локальный перенос узла между соседними группами проходов;
 \item конечная энергия, устойчивость, спин и статистика.
\end{enumerate}

\begin{nogocriterion}[Запрет переименования эффективного вихря]
Нельзя объявлять построенную вихревую нить фундаментальной нитью ткани.
Она остаётся эффективным полевым дефектом, пока не выведено отображение от
одной микроскопической связной линии к полям $Q$, $T$ и $B$.
\end{nogocriterion}

\section{Первый различающий тест}

Нулевая гипотеза: одно и то же неотрицательное локальное распределение
на четырёх тетраэдрических направлениях одновременно воспроизводит
упорядоченный второй момент $Q_*$, канонический третий момент $T_*$ и
нулевой полярный ток замкнутой нити.

Стоп-критерий: если эти три условия несовместимы, гипотеза одной нити не
закрывается, но поля $Q$ и $T$ запрещено считать обычными моментами одной
и той же локальной плотности. Тогда потребуется фазовое или
связностно-взвешенное чтение третьего момента.

\begin{protocol}
\begin{enumerate}
 \item спокойное одиночное кольцо: \textbf{закрыто отрицательно};
 \item вихрь как эффективное ребро ткани: \textbf{сохранён};
 \item одна глобальная нить: \textbf{гипотеза, не доказательство};
 \item ненулевое микроскопическое сечение: \textbf{требует вывода};
 \item планковская нормировка сечения: \textbf{не выведена};
 \item новый первичный тест: \textbf{совместная моментная реализация};
 \item кривизный ответ одиночной трубки: \textbf{вторичный маршрут};
 \item рождение материи: \textbf{не закрыто}.
\end{enumerate}
\end{protocol}